\documentclass{article}
\usepackage{graphicx} % Required for inserting images
\usepackage[russian]{babel}
\usepackage{graphicx}
\usepackage{hyperref}
\title{Мини-конспект по теме: Теорема Пифагора}
\author{Автор: Щетнев Алексей}
\date{16 июля 2025 г}

\begin{document}


\maketitle

\tableofcontents

\newpage
\section{Введение}
Теорема Пифагора - одна из важнейших теорем евклидной геометрии. Она находит
применение в самых разных областях:

\begin{itemize}
    \item геометрия и тригонометрия
    \item физика
    \item инженерные науки
    \item компьютерная графика
\end{itemize}
 
\section{Формулировка теоремы}

\textbf{Слова:} В прямоугольном треугольнике квадрат гипотенузы равен сумме квадратов катетов.
\begin{equation}
c^2 = a^2 + b^2
\end{equation}
Как видно из формулы 1, значение двух сторон позволяет найти третью

\section{Доказательство (набросок)}
Одно из доказательст основывается на площади квадрата, составленного из четырёх одинкаовых прямоугольников и малого квадрата в центре. Раскладывая площадь двумя способами, получаем c^2 = a^2 + b^2

\section{Примеры расчёта}

\textbf{Пример 1}
\[
a=3, b=4
\]
\[
c=\sqrt{a^2+b^2}=\sqrt{9+16}=5
\]
\textbf{Пример 2}
\begin{enumerate}
    \item Дано а=5,b=12
    \item Решение:
\end{enumerate}
\[
c=\sqrt{5^2+12^2 }=\sqrt{25+144}=13
\]


\section{Таблица значений}
\begin{center}
    
\begin{tabular}{|c|c|c|}
\hline     
Катет a & Катет b & Гипотенуза c \\
\hline
3 & 4 & 5 \\
\hline
5 & 12 & 13 \\
\hline
7&24&25\\
\hline
\end{tabular}

\end{center}

\newpage

\section{Иллюстрация}

Ниже пример изображения (не забудьте добавить файл triangle.png в ту же папку):

\includegraphics[width=0.5\textwidth]{ц66.png}

\section{Заключение}

Теорема Пифагора - один из краеугольных камней геометрии, помогающий решать множество практических задач.

\section{Ссылки и литература}

\begin{itemize}
    \item \href{https://ru.wikipedia.org/wiki/%D0%A2%D0%B5%D0%BE%D1%80%D0%B5%D0%BC%D0%B0_%D0%9F%D0%B8%D1%84%D0%B0%D0%B3%D0%BE%D1%80%D0%B0}{Википедия: Теорема Пифагора}
    \item \href{https://prosv.ru/catalog/geometriya56-8-kl/?ysclid=mupbjw3ce9631886878}{Классические учебники геометрии}
\end{itemize}

\end{document}
